\chapter{Struktury, wartościowania i prawda}\label{ch:semantics-fol}

\section{Co oznacza nadanie znaczenia napisowi?}
Ten sam napis $R(v_0,v_1)$ może oznaczać „$v_0$ jest mniejszy od $v_1$”,
„$v_0$ i $v_1$ są sąsiadami w grafie” lub zupełnie inną relację.
Składnia z rozdziału~\ref{ch:syntax-fol} nie rozstrzyga, która z interpretacji
jest właściwa. Znaczenie przypisujemy dopiero przez wskazanie zbioru obiektów
i interpretacji symboli pozalogicznych. Następnie definiujemy, kiedy
konkretna formuła jest spełniona. Ta druga definicja jest indukcyjna:
wartości formuł złożonych zależą od wartości ich części.

\section{Struktury i ich uniwersa}
\begin{definicja}[Struktura dla sygnatury]\label{def:sem-structure}
Niech $\mathcal L$ będzie sygnaturą z Definicji~\ref{def:syntax-signature}.
\emph{Strukturą} $\mathfrak A$ dla $\mathcal L$ nazywamy niepusty zbiór $A$,
zwany jej \emph{uniwersum}, razem z następującymi danymi:
\begin{enumerate}[label=(\roman*)]
\item dla każdego $c\in\operatorname{Const}(\mathcal L)$ element
      $c^{\mathfrak A}\in A$;
\item dla każdego symbolu funkcyjnego $f$ arności $n$ funkcja
      $f^{\mathfrak A}:A^n\to A$;
\item dla każdego symbolu relacyjnego $R$ arności $n$ relacja
      $R^{\mathfrak A}\subseteq A^n$.
\end{enumerate}
Symbol równości $=$ jest zawsze interpretowany jako rzeczywista równość
w $A$; nie jest elementem sygnatury pozalogicznej.
\end{definicja}
W tym skrypcie przyjmujemy standardową konwencję, że uniwersum jest
niepuste. Nie ogranicza to rozważań o pustych zbiorach jako obiektach:
na przykład uniwersum modelu teorii mnogości może zawierać zbiór pusty,
ale samo uniwersum modelu nie jest puste.

\begin{przyklad}[Dwie interpretacje tej samej sygnatury]
Weźmy język z jednym binarnym symbolem $R$.
Niech $\mathfrak A$ ma uniwersum $\{0,1\}$ i interpretację
$R^{\mathfrak A}=\{(0,1)\}$.
Niech $\mathfrak B$ ma to samo uniwersum, lecz
$R^{\mathfrak B}=\{(0,0),(1,1)\}$.
Wówczas zdanie $\exists v_0\,R(v_0,v_0)$ jest fałszywe
w $\mathfrak A$, a prawdziwe w $\mathfrak B$.
Mamy tę samą składnię, ale różne interpretacje.
\end{przyklad}
\begin{przyklad}[Liczby naturalne i grafy]
W języku arytmetyki standardową strukturą jest
$\mathfrak N=(\N,0,S,+,\cdot,<)$, gdzie $S(n)=n+1$.
W języku grafów struktura ma postać $(V,E)$, gdzie $V\ne\varnothing$
i $E\subseteq V\times V$. Formuła $\forall v_0\,\neg E(v_0,v_0)$
wyraża brak pętli: żaden wierzchołek nie jest połączony ze sobą.
\end{przyklad}


\begin{uwaga}[Stałe, funkcje i relacje pełnią inne role]
W strukturze arytmetycznej symbol $0$ wskazuje pojedynczy element,
$S$ i $+$ wskazują funkcje przypisujące wyniki argumentom,
a $<$ wskazuje relację, której pary mogą być prawdziwe albo fałszywe.
Dlatego napis $S(0)$ jest termem (oznacza obiekt), lecz $0<S(0)$
jest formułą (wyraża własność dwóch obiektów).
Gdy w następnym kroku przypiszemy zmiennej $v_0$ wartość $2$,
term $S(v_0)$ będzie oznaczał $3$, a formuła $v_0<S(v_0)$
będzie prawdziwa. Są to dwa różne rodzaje obliczeń.
\end{uwaga}

\section{Wartościowania i interpretacja termów}
\begin{definicja}[Wartościowanie]\label{def:sem-assignment}
Niech $\mathfrak A$ będzie strukturą o uniwersum $A$.
\emph{Wartościowaniem zmiennych} nazywamy funkcję
$s:\{v_0,v_1,\ldots\}\to A$.
Dla $a\in A$ i zmiennej $x$ symbol $s[x\mapsto a]$
oznacza wartościowanie przypisujące $x$ wartość $a$ i zgadzające
się z $s$ na wszystkich pozostałych zmiennych.
\end{definicja}
\begin{uwaga}[Dlaczego wartościowania istnieją?]
Przyjmujemy $A\ne\varnothing$. Po wybraniu jednego
elementu $a_0\in A$ możemy przypisać każdej zmiennej
wartość $a_0$; otrzymujemy wartościowanie stałe.
Nie wymaga to wyboru nieskończonej rodziny elementów:
jedno ustalone $a_0$ wystarcza dla wszystkich zmiennych.
Przy pustym uniwersum takiego wartościowania dla
nieskończonego zbioru zmiennych nie byłoby.
\end{uwaga}

\begin{definicja}[Wartość termu]\label{def:sem-term-value}
Wartość termu $t$ w strukturze $\mathfrak A$ względem wartościowania $s$
oznaczamy przez $t^{\mathfrak A}[s]$ i określamy indukcyjnie:
\begin{enumerate}[label=(\roman*)]
\item $x^{\mathfrak A}[s]=s(x)$ dla zmiennej $x$;
\item $c^{\mathfrak A}[s]=c^{\mathfrak A}$ dla stałej $c$;
\item jeśli $t=f(t_1,\ldots,t_n)$, to
      $t^{\mathfrak A}[s]=f^{\mathfrak A}(t_1^{\mathfrak A}[s],\ldots,t_n^{\mathfrak A}[s])$.
\end{enumerate}
\end{definicja}
\begin{przyklad}
Przy $s(v_0)=3$, $s(v_1)=5$ w $\mathfrak N$ term
$t=+(S(v_0),v_1)$ ma wartość
$t^{\mathfrak N}[s]=S^{\mathfrak N}(3)+5=9$.
Term domknięty $S(S(0))$ ma wartość $2$ niezależnie od $s$.
\end{przyklad}

\begin{lemat}[Zależność wartości termu od jego zmiennych]\label{lem:sem-term-coincidence}
Niech $t$ będzie termem, a $s,r$ wartościowaniami w tej samej
strukturze $\mathfrak A$. Jeżeli $s(x)=r(x)$ dla każdej zmiennej
$x\in\operatorname{FV}(t)$, to
$t^{\mathfrak A}[s]=t^{\mathfrak A}[r]$.
W szczególności wartość termu domkniętego nie zależy od wartościowania.
\end{lemat}
\begin{proof}
Przeprowadzimy indukcję po konstrukcji termu. Jeśli $t=x$,
założenie daje $s(x)=r(x)$, a więc równość wartości. Jeśli $t=c$
jest stałą, obie wartości są równe $c^{\mathfrak A}$.
Jeśli $t=f(t_1,\ldots,t_n)$, to z definicji zbioru zmiennych
wolnych $\operatorname{FV}(t_i)\subseteq\operatorname{FV}(t)$.
Zatem $s$ i $r$ zgadzają się na zmiennych każdego $t_i$.
Z założenia indukcyjnego dostajemy
$t_i^{\mathfrak A}[s]=t_i^{\mathfrak A}[r]$ dla $i=1,\ldots,n$.
Przyłożenie tej samej funkcji $f^{\mathfrak A}$ do identycznych
$n$-tek daje równość wartości $t$.
Jeśli $t$ jest domknięty, zbiór $\operatorname{FV}(t)$ jest pusty,
więc warunek zgodności wartościowań jest spełniony automatycznie.
\end{proof}

\section{Relacja spełniania: definicja rekurencyjna}
\begin{definicja}[Spełnianie formuły]\label{def:sem-satisfaction}
Niech $\mathcal L$ będzie sygnaturą, $\mathfrak A$ strukturą
języka $\mathcal L$ o uniwersum $A$, a $s$ wartościowaniem.
Dla każdej formuły $\varphi$ języka $\mathcal L$ napis
$\mathfrak A,s\models\varphi$ oznacza, że $\mathfrak A$ spełnia
formułę $\varphi$ przy $s$. Definiujemy tę relację indukcyjnie:
\begin{enumerate}[label=(\roman*)]
\item $\mathfrak A,s\models t_1=t_2$ wtedy i tylko wtedy, gdy
      $t_1^{\mathfrak A}[s]=t_2^{\mathfrak A}[s]$;
\item $\mathfrak A,s\models R(t_1,\ldots,t_n)$ wtedy i tylko wtedy, gdy
      $(t_1^{\mathfrak A}[s],\ldots,t_n^{\mathfrak A}[s])\in R^{\mathfrak A}$;
\item $\mathfrak A,s\models\neg\varphi$ wtedy i tylko wtedy, gdy
      nie zachodzi $\mathfrak A,s\models\varphi$;
\item $\mathfrak A,s\models\varphi\land\psi$ wtedy i tylko wtedy,
      gdy spełnione są obie formuły; dla $\lor$, $\to$ i $\leftrightarrow$
      używamy tablic prawdy z rozdziału~\ref{ch:propositional};
\item $\mathfrak A,s\models\forall x\,\varphi$ wtedy i tylko wtedy,
      gdy dla każdego $a\in A$ zachodzi
      $\mathfrak A,s[x\mapsto a]\models\varphi$;
\item $\mathfrak A,s\models\exists x\,\varphi$ wtedy i tylko wtedy,
      gdy dla pewnego $a\in A$ zachodzi
      $\mathfrak A,s[x\mapsto a]\models\varphi$.
\end{enumerate}
\end{definicja}
W szczególności formuła $\varphi\to\psi$ jest fałszywa dokładnie wtedy,
gdy $\varphi$ jest prawdziwa i $\psi$ fałszywa przy tym samym wartościowaniu.
Dla kwantyfikatorów zmieniamy wartość wyłącznie związanej zmiennej.

\begin{stwierdzenie}[Jednoznaczność spełniania]
\label{prop:sem-uniqueness-xlv}
Dla każdej sygnatury $\mathcal L$, struktury $\mathfrak A$
o niepustym uniwersum i wartościowania $s$ reguły
Definicji~\ref{def:sem-satisfaction} wyznaczają jednoznacznie,
które formuły $\mathcal L$ są spełnione przy $s$.
\end{stwierdzenie}
\begin{proof}
Najpierw Lematu~\ref{lem:syntax-unique-terms} używamy
rekurencyjnie do wyznaczenia wartości wszystkich termów:
zmienne i stałe mają jedną wskazaną wartość, a wartość
aplikacji $f$ jest zadana przez funkcję $f^{\mathfrak A}$.
Z Lematu~\ref{lem:syntax-unique-formulas} wynika, że
każda formuła ma jeden z określonych kroków końcowych.
Dla atomu reguły równości albo przynależności do relacji
wyznaczają wartość logiczną. Dla negacji i spójników
binarnych wartość wynika jednoznacznie z wartości
podformuł. Dla kwantyfikatora $\forall x$ rozstrzygamy,
czy wskazana podformuła jest spełniona przy
$s[x\mapsto a]$ dla wszystkich $a\in A$; dla $\exists x$
sprawdzamy, czy zachodzi to dla pewnego $a\in A$.
Indukcja po skończonej budowie formuły dowodzi zatem,
że żaden jej krok nie dopuszcza dwóch różnych ocen.
\end{proof}

\begin{przyklad}[Parzystość jako formuła]
Niech $\varphi(v_0)$ będzie formułą
$\exists v_1\,(+(v_1,v_1)=v_0)$.
W modelu $\mathfrak N$ przy wartościowaniu z $s(v_0)=8$
formuła jest spełniona: wybieramy wartość $4$ dla $v_1$.
Przy wartościowaniu $r(v_0)=7$ nie jest spełniona,
bo $n+n=7$ nie ma rozwiązania w $\N$.
Nie pytamy tu o „prawdę formuły bez kontekstu”: wynik zależy
od wartości zmiennej wolnej $v_0$.
\end{przyklad}

\begin{przyklad}[Kolejność kwantyfikatorów]
W języku z relacją $R$ rozważmy strukturę z uniwersum $\N$,
w której $R^{\mathfrak A}$ oznacza zwykły porządek $<$.
Wtedy $\forall x\,\exists y\,R(x,y)$ jest prawdziwe:
dla danego $n$ wystarczy wziąć $n+1$.
Natomiast $\exists y\,\forall x\,R(x,y)$ jest fałszywe:
dla dowolnego kandydata $m$ można wziąć $x=m$ i otrzymać
fałszywe $m<m$. Te dwa napisy różnią się kolejnością kwantyfikatorów
— i naprawdę wyrażają odmienne twierdzenia.
\end{przyklad}

\begin{przyklad}[Obliczenie spełniania w skończonej strukturze]
Niech $A=\{0,1\}$, a $R^{\mathfrak A}=\{(0,0),(1,1)\}$.
Dla każdego $a\in A$ istnieje dokładnie jeden $b\in A$
z własnością $(a,b)\in R^{\mathfrak A}$: jest nim $b=a$.
Stąd $\mathfrak A\models\forall x\,\exists y\,R(x,y)$.
Nie istnieje jednak jeden element $b\in A$, dla którego
jednocześnie $(0,b)$ i $(1,b)$ należałyby do tej relacji.
Zatem $\mathfrak A\not\models\exists y\,\forall x\,R(x,y)$.
Rozpisując świadków, dla $x=0$ wybieramy $y=0$,
a dla $x=1$ wybieramy $y=1$. W drugim zdaniu $y$ musiałoby
być wybrane \emph{przed} $x$, więc te dwa wybory nie wystarczą.
\end{przyklad}


\begin{twierdzenie}[Prawa kwantyfikatorów w semantyce pierwszego rzędu]
\label{thm:sem-quantifier-laws-xlv}
Niech $\mathcal L$ będzie sygnaturą, $\mathfrak A$ jej
strukturą o niepustym uniwersum $A$, $s$ wartościowaniem,
$\varphi$ formułą, a $x,y$ zmiennymi. Wtedy:
\begin{align*}
\mathfrak A,s\models\neg\forall x\,\varphi
 &\iff\mathfrak A,s\models\exists x\,\neg\varphi,\\
\mathfrak A,s\models\neg\exists x\,\varphi
 &\iff\mathfrak A,s\models\forall x\,\neg\varphi,\\
\mathfrak A,s\models\forall x\,\forall y\,\varphi
 &\iff\mathfrak A,s\models\forall y\,\forall x\,\varphi,\\
\mathfrak A,s\models\exists x\,\exists y\,\varphi
 &\iff\mathfrak A,s\models\exists y\,\exists x\,\varphi.
\end{align*}
Nie twierdzimy, że można dowolnie zamieniać
kwantyfikator ogólny z egzystencjalnym.
\end{twierdzenie}
\begin{proof}
Pierwsza formuła po lewej oznacza, że nie jest prawdą,
iż dla każdego $a\in A$ zachodzi
$\mathfrak A,s[x\mapsto a]\models\varphi$.
Ponieważ negujemy kwantyfikator ogólny, istnieje
$a\in A$ takie, że
$\mathfrak A,s[x\mapsto a]\not\models\varphi$.
Z definicji negacji jest to dokładnie warunek
$\mathfrak A,s[x\mapsto a]\models\neg\varphi$,
czyli prawdziwość $\exists x\,\neg\varphi$ przy $s$.
Odwrotny kierunek czytamy w tej samej równoważności
od końca do początku.

Podobnie zaprzeczenie $\exists x\,\varphi$ oznacza,
że dla żadnego $a\in A$ formuła $\varphi$ nie jest
spełniona przy $s[x\mapsto a]$. To znaczy, że
$\neg\varphi$ jest spełniona przy każdym takim
wartościowaniu, co jest definicją $\forall x\,\neg\varphi$.

Rozpiszmy trzecią równoważność: lewa strona wymaga,
aby dla każdego $a\in A$ i każdego $b\in A$ zachodziło
$\mathfrak A,s[x\mapsto a][y\mapsto b]\models\varphi$.
Gdy $x\ne y$, operacje przypisania wartości różnym
zmiennym komutują, więc otrzymujemy dokładnie warunek
po prawej, z $b$ wybranym przed $a$. Gdy $x=y$,
obydwa zdania są identyczne i równoważność jest oczywista.
Dla dwóch kwantyfikatorów egzystencjalnych zamiast
warunku dla wszystkich par $(a,b)$ otrzymujemy istnienie
przynajmniej jednej pary; kolejność jej współrzędnych
nie wpływa na istnienie świadka. Tak samo obsługujemy
przypadek $x=y$.
\end{proof}
\begin{przyklad}[Zakres prawa de Morgana dla kwantyfikatorów]
Niech $A=\{0,1\}$, a $P^{\mathfrak A}=\{0\}$ dla
jednoargumentowego symbolu relacyjnego $P$.
Zdanie $\forall x\,P(x)$ jest fałszywe, gdyż
$1\notin P^{\mathfrak A}$. Jego negacja jest więc
prawdziwa. Zarazem $1$ świadczy o prawdziwości
$\exists x\,\neg P(x)$.
Nie wynika natomiast równoważność
$\forall x\,\exists y\,R(x,y)$ i
$\exists y\,\forall x\,R(x,y)$; kontrprzykład
ze skończonej struktury podano bezpośrednio wyżej.
\end{przyklad}

\section{Lemat o zgodności wartościowań}
\begin{twierdzenie}[Spełnienie zależy tylko od zmiennych wolnych]
\label{thm:sem-coincidence}
Niech $\mathcal L$ będzie sygnaturą, $\varphi$ formułą
języka $\mathcal L$, $\mathfrak A$ strukturą tego języka
o niepustym uniwersum $A$, a $s,r$ jej wartościowaniami
zgadzającymi się na
$\operatorname{FV}(\varphi)$. Wtedy
\[
\mathfrak A,s\models\varphi
\quad\Longleftrightarrow\quad
\mathfrak A,r\models\varphi.
\]
\end{twierdzenie}
\begin{proof}
Dowód przeprowadzamy indukcją po budowie $\varphi$.
Dla formuł atomowych wystąpienia zmiennych wolnych pochodzą
z termów. Zgodność wartościowań na tych zmiennych oraz
Lemat~\ref{lem:sem-term-coincidence} dają identyczne wartości
wszystkich termów. Stąd równoważność spełniania obu typów atomów.

Dla negacji teza przechodzi natychmiast: jeśli $\psi$ jest
prawdziwa przy obu wartościowaniach albo fałszywa przy obu,
to tak samo jest z $\neg\psi$. Dla każdego spójnika binarnego
stosujemy założenie indukcyjne do obu podformuł, ponieważ
ich zmienne wolne należą do $\operatorname{FV}(\varphi)$;
tablica prawdy daje ten sam wynik.

Pozostał przypadek $\varphi=\forall x\,\psi$.
Załóżmy, że $s$ i $r$ zgadzają się na
$\operatorname{FV}(\psi)\setminus\{x\}$.
Dla dowolnego $a\in A$ wartościowania
$s[x\mapsto a]$ i $r[x\mapsto a]$ zgadzają się na
wszystkich zmiennych wolnych $\psi$: na $x$ mają tę samą
wartość $a$, a na pozostałych wymaganą zgodność zapewnia założenie.
Z założenia indukcyjnego wynika więc
\[
\mathfrak A,s[x\mapsto a]\models\psi
\quad\Longleftrightarrow\quad
\mathfrak A,r[x\mapsto a]\models\psi.
\]
Ponieważ zachodzi to dla \emph{każdego} $a\in A$,
definicja $\forall$ daje tezę. W przypadku
$\exists x\,\psi$ ten sam argument stosujemy do istnienia
odpowiedniej wartości $a$, zamiast do wszystkich $a$.
\end{proof}

\begin{wniosek}[Prawda zdania jest niezależna od wartościowania]
\label{cor:sem-sentence}
Jeżeli $\sigma$ jest zdaniem i $\mathfrak A$ strukturą,
to albo $\mathfrak A,s\models\sigma$ dla wszystkich wartościowań $s$,
albo dla żadnego. Możemy więc pisać $\mathfrak A\models\sigma$.
\end{wniosek}
\begin{proof}
Zdanie nie ma zmiennych wolnych, zatem każde dwa wartościowania
spełniają założenie Twierdzenia~\ref{thm:sem-coincidence}.
\end{proof}

\section{Podstawienie termu a jego wartość}
Lemat z tego paragrafu tłumaczy, dlaczego reguły podstawiania w dowodach
muszą zabraniać przechwytywania zmiennych. Bez tego formuła
po podstawieniu mogłaby wyrażać coś innego niż zamierzona instancja.

W zapisie $t[x:=u]$ stosujemy
Definicję~\ref{def:syntax-term-substitution-xlv}; w szczególności
podstawienie termu nie wymaga warunku unikania przechwycenia,
bo termy nie zawierają kwantyfikatorów.

\begin{lemat}[Podstawienie w termie]\label{lem:sem-term-subst}
Niech $u,t$ będą termami, $x$ zmienną, $\mathfrak A$ strukturą
i $s$ wartościowaniem. Wówczas
\[
(t[x:=u])^{\mathfrak A}[s]
=t^{\mathfrak A}[s[x\mapsto u^{\mathfrak A}[s]]].
\]
\end{lemat}
\begin{proof}
Prowadzimy indukcję po $t$. Jeżeli $t=x$, lewa strona jest
$u^{\mathfrak A}[s]$, a prawa wartością $x$ przy wartościowaniu
przypisującym $x$ dokładnie ten sam element uniwersum $A$.
Dla innej zmiennej $y$ podstawienie niczego nie zmienia,
a oba wartościowania przypisują $y$ tę samą wartość.
Dla stałej obie strony są jej interpretacją.
Jeżeli $t=f(t_1,\ldots,t_n)$, stosujemy założenie indukcyjne
do każdego $t_i$, a następnie do obu otrzymanych $n$-tek
przykładamy funkcję $f^{\mathfrak A}$.
\end{proof}

\begin{twierdzenie}[Lemat o podstawieniu]\label{thm:sem-substitution}
Niech $\varphi$ będzie formułą, $x$ zmienną, a $t$ termem
dopuszczalnym do podstawienia za $x$ w $\varphi$.
Dla każdej struktury $\mathfrak A$ i wartościowania $s$ zachodzi
\[
\mathfrak A,s\models\varphi[x:=t]
\quad\Longleftrightarrow\quad
\mathfrak A,s[x\mapsto t^{\mathfrak A}[s]]\models\varphi.
\]
\end{twierdzenie}
\begin{proof}
Prowadzimy indukcję po budowie $\varphi$. Najpierw
rozpatrzmy atom równościowy $u_1=u_2$. Po lewej
podstawiamy term $t$ do obu argumentów, więc zgodnie
z Lematem~\ref{lem:sem-term-subst} ich wartości przy $s$
równe są wartościom $u_1,u_2$ przy
$s[x\mapsto t^{\mathfrak A}[s]]$. Definicja równości
atomowej daje żądaną równoważność. Dla atomu relacyjnego
wszystkie argumenty mają odpowiednio te same wartości,
więc należą lub nie należą do tej samej relacji
$R^{\mathfrak A}$. Dla negacji i każdego spójnika
binarnego stosujemy założenie indukcyjne do odpowiednich
podformuł i reguły spełniania z
Definicji~\ref{def:sem-satisfaction}.

Pozostały kwantyfikatory. Niech $\varphi=Qy\,\psi$,
gdzie $Q\in\{\forall,\exists\}$. Rozważamy trzy przypadki.
\begin{enumerate}[label=(\roman*)]
\item Gdy $y=x$, kwantyfikator wiąże wszystkie wystąpienia
$x$ wolne w $\psi$. Stąd $x\notin\operatorname{FV}(\varphi)$,
a $\varphi[x:=t]=\varphi$. Wartościowania
$s$ i $s[x\mapsto t^{\mathfrak A}[s]]$ zgadzają się
na zmiennych wolnych $\varphi$; zgodnie z
Twierdzeniem~\ref{thm:sem-coincidence} spełniają ją
jednocześnie albo oba nie spełniają.
\item Gdy $y\ne x$, lecz $x\notin\operatorname{FV}(\psi)$,
podstawienie nic nie zmienia. Ponieważ wtedy również
$x\notin\operatorname{FV}(\varphi)$, tezę ponownie
daje twierdzenie o zgodności wartościowań.
\item Gdy $y\ne x$ oraz $x\in\operatorname{FV}(\psi)$,
każde wolne wystąpienie $x$ w $\psi$ pozostaje w zakresie
$Qy$. Dopuszczalność podstawienia implikuje więc
$y\notin\operatorname{FV}(t)$. Dla dowolnego $a\in A$
założenie indukcyjne zastosowane do $\psi$ i
$s[y\mapsto a]$ daje
\[
\mathfrak A,s[y\mapsto a]\models\psi[x:=t]
\iff
\mathfrak A,s[y\mapsto a]
  [x\mapsto t^{\mathfrak A}[s[y\mapsto a]]]
  \models\psi.
\]
Lemat~\ref{lem:sem-term-coincidence} pozwala zastąpić
$t^{\mathfrak A}[s[y\mapsto a]]$ przez
$t^{\mathfrak A}[s]$, gdyż $y$ nie jest zmienną termu $t$.
Dla różnych zmiennych $x,y$ operacje aktualizacji
wartościowania komutują, zatem ostatnie wartościowanie to
$s[x\mapsto t^{\mathfrak A}[s]][y\mapsto a]$.
Równoważność zachodzi dla każdego $a\in A$.
Jeśli $Q=\forall$, zastosowanie warunku \emph{dla wszystkich}
$a$ daje tezę; jeśli $Q=\exists$, warunku
\emph{dla pewnego} $a$.
\end{enumerate}
Każdy przypadek tworzenia formuł został rozpatrzony.
\end{proof}
\begin{przyklad}[Zmiana nazwy zmiennej związanej]
Formuły $\forall v_0\,(v_0=v_1)$ i
$\forall v_2\,(v_2=v_1)$ mają takie samo znaczenie w każdej
strukturze przy każdym wartościowaniu: w obu przypadkach
sprawdzamy równość wszystkich elementów uniwersum z wartością $v_1$.
Nie wolno natomiast dowolnie przemianowywać zmiennych tak,
aby zmienić zakres kwantyfikatora lub przechwycić wolne wystąpienie.
\end{przyklad}

\begin{przyklad}[Dlaczego lemat o podstawieniu wymaga dopuszczalności]
Niech $A=\{0,1\}$, $\varphi=\forall y\,(x=y)$,
$t=y$ oraz $s(x)=0$, $s(y)=1$. Naiwna zamiana wolnego
$x$ na $y$ daje $\forall y\,(y=y)$, prawdziwą w $A$.
Natomiast po prawej stronie wzoru z
Twierdzenia~\ref{thm:sem-substitution} występuje
$\varphi$ przy $s[x\mapsto t^{\mathfrak A}[s]]$,
czyli przy $x=1$. Formuła $\forall y\,(1=y)$ jest
fałszywa, ponieważ dla $y=0$ otrzymujemy $1\ne0$.
Wobec tego równoważność z lematu o podstawieniu bez
warunku dopuszczalności jest fałszywa. Po przemianowaniu
zmiennej związanej na świeżą $z$ otrzymujemy poprawne
$\forall z\,(y=z)$, fałszywe przy $s(y)=1$.
\end{przyklad}

\begin{twierdzenie}[Poprawność przemianowania zmiennej związanej]
\label{thm:sem-alpha-renaming-xlv}
Niech $\mathfrak A$ będzie strukturą języka $\mathcal L$,
$\psi$ formułą tego języka, a $x,y$ dwiema różnymi
zmiennymi. Załóżmy, że $y$ \emph{nie występuje w ogóle}
w napisie $\psi$ (ani wolno, ani związanie).
Wtedy dla każdego wartościowania $s$ oraz
$Q\in\{\forall,\exists\}$ zachodzi
\[
\mathfrak A,s\models Qx\,\psi
\quad\Longleftrightarrow\quad
\mathfrak A,s\models Qy\,\psi[x:=y].
\]
Podstawienie po prawej jest dopuszczalne.
\end{twierdzenie}
\begin{proof}
Ponieważ $y$ nie występuje w $\psi$, żaden kwantyfikator
w jej wnętrzu nie może przechwycić $y$ podstawianego za $x$.
Dla dowolnego $a\in A$ stosujemy
Twierdzenie~\ref{thm:sem-substitution} do formuły $\psi$,
termu $y$ i wartościowania $s[y\mapsto a]$:
\[
\mathfrak A,s[y\mapsto a]\models\psi[x:=y]
\iff
\mathfrak A,s[y\mapsto a][x\mapsto a]\models\psi.
\]
Wartościowania po prawej oraz $s[x\mapsto a]$ różnią się
co najwyżej na zmiennej $y$. Ponieważ
$y\notin\operatorname{FV}(\psi)$, twierdzenie o zgodności
wartościowań daje
\[
\mathfrak A,s[y\mapsto a]\models\psi[x:=y]
\iff
\mathfrak A,s[x\mapsto a]\models\psi.
\]
Dla $Q=\forall$ wymagamy tej równoważności dla każdego
$a\in A$, a dla $Q=\exists$ wystarczy odpowiedni
świadek $a$. Otrzymujemy oba kierunki w każdym przypadku.
\end{proof}
\begin{przyklad}[Kiedy przemianowanie jest błędne?]
Niech $\psi$ ma postać $x=y$. Wtedy $y$ nie jest świeża.
Nie wolno zastąpić $\forall x\,(x=y)$ przez
$\forall y\,(y=y)$: w strukturze
$A=\{0,1\}$ pierwsza formuła jest fałszywa przy każdym
wartościowaniu, natomiast druga jest prawdziwa.
Jeżeli zamiast $y$ weźmiemy świeżą zmienną $z$,
formuła $\forall z\,(z=y)$ jest równoważna wyjściowej
przy każdym wartościowaniu, zgodnie z twierdzeniem.
\end{przyklad}

\section{Modele, kontrmodele i konsekwencja semantyczna}
\begin{definicja}[Model i konsekwencja semantyczna]\label{def:sem-entailment}
Niech $\mathcal L$ będzie sygnaturą; wszystkie rozważane
zdania i struktury należą do tego samego języka $\mathcal L$.
Jeżeli $\sigma$ jest zdaniem, to $\mathfrak A\models\sigma$
oznacza jego prawdziwość w $\mathfrak A$.
Dla zbioru zdań $\Gamma$ piszemy $\mathfrak A\models\Gamma$,
jeśli $\mathfrak A$ spełnia każde zdanie z $\Gamma$;
wówczas $\mathfrak A$ nazywamy \emph{modelem} $\Gamma$.
Dla zdania $\sigma$ zapis $\Gamma\models\sigma$ oznacza, że
każdy model $\Gamma$ jest również modelem $\sigma$.
Jeżeli $\Gamma=\varnothing$, piszemy $\models\sigma$
i nazywamy $\sigma$ \emph{logicznie prawdziwym} (ważnym).
\emph{Kontrmodel} dla $\Gamma\models\sigma$ to struktura,
która spełnia $\Gamma$, lecz nie spełnia $\sigma$.
\end{definicja}
W tym rozdziale wystarczają zbiory \emph{zdań}. W dalszej części
skryptu rozszerzymy tę relację na formuły ze zmiennymi
wolnymi, zachowując zgodność z powyższą definicją.

\begin{przyklad}[Kontrmodel dla błędnej zmiany kwantyfikatorów]
Rozważmy formułę, której głównym spójnikiem jest implikacja
\[
\bigl(\forall x\,\exists y\,R(x,y)\bigr)
\to\bigl(\exists y\,\forall x\,R(x,y)\bigr).
\]
Nie jest ona logicznie prawdziwa. W strukturze o uniwersum $\N$,
w której $R$ oznacza $<$, poprzednik jest prawdziwy (można wziąć $y=x+1$),
a następnik fałszywy (dla dowolnego $y$ wybieramy $x=y$).
Jedna taka struktura wystarcza, aby obalić powszechną ważność implikacji.
\end{przyklad}
\begin{przyklad}[Ważność a prawdziwość w jednym modelu]
Zdanie $\forall x\,(x=x)$ jest logicznie prawdziwe:
każdy element jest równy samemu sobie, niezależnie od struktury.
Zdanie $\exists x\,\exists y\,\neg(x=y)$ jest prawdziwe
w strukturze o co najmniej dwóch elementach, ale fałszywe
w strukturze jednoelementowej. Nie jest więc logicznie prawdziwe.
\end{przyklad}

\begin{twierdzenie}[Monotoniczność konsekwencji]
\label{thm:sem-entailment-monotone-xlv}
Dla dowolnych zbiorów zdań $\Gamma\subseteq\Delta$
tej samej sygnatury i dowolnego zdania $\sigma$ zachodzi:
$\Gamma\models\sigma$ pociąga $\Delta\models\sigma$.
Jeżeli $\Gamma$ nie ma modelu, to $\Gamma\models\sigma$
dla każdego zdania $\sigma$ tej sygnatury.
\end{twierdzenie}
\begin{proof}
Niech $\mathfrak A$ będzie dowolnym modelem $\Delta$.
Ponieważ każdy element $\Gamma$ należy do $\Delta$,
struktura $\mathfrak A$ spełnia wszystkie zdania $\Gamma$.
Z $\Gamma\models\sigma$ otrzymujemy
$\mathfrak A\models\sigma$. Wybór modelu $\Delta$
był dowolny, więc $\Delta\models\sigma$.
Jeżeli $\Gamma$ nie ma modeli, nie istnieje kontrmodel:
żadna struktura nie spełnia wszystkich przesłanek
$\Gamma$ i jednocześnie nie spełnia $\sigma$.
Warunek definiujący $\Gamma\models\sigma$ zachodzi
wówczas dla każdego $\sigma$.
\end{proof}
\begin{przyklad}[Konsekwencja z niespełnialnych przesłanek]
Niech $\Gamma=\{\forall x\,\neg(x=x)\}$.
Każde uniwersum jest niepuste, więc wybieramy $a\in A$.
Atom $a=a$ jest prawdziwy, stąd
$\forall x\,\neg(x=x)$ jest fałszywe w każdej strukturze.
Zatem $\Gamma$ nie ma modelu i
$\Gamma\models\exists x\,\neg(x=x)$, chociaż zdanie
$\exists x\,\neg(x=x)$ nie jest logicznie prawdziwe.
Nie należy mylić \emph{wynikania z przesłanek}
z \emph{ważnością bez przesłanek}.
\end{przyklad}

\section{Rozszerzenie języka: logika drugiego rzędu}\label{sec:second-order-introduction}
\subsection{Jakie pytania są niedostępne w pierwszym rzędzie?}
W logice pierwszego rzędu formuła $\forall x\,\varphi(x)$ mówi o
\emph{wszystkich elementach} uniwersum $A$. Nawet jeśli dla dowolnie
wybranego podzbioru $D\subseteq A$ potrafimy opisać, że pewne zdanie
dotyczy $D$, sam kwantyfikator $\forall x$ nie obejmuje zbiorów $D$.
Aby móc powiedzieć na przykład „każdy zbiór zawierający zero
i zamknięty na następnik zawiera wszystkie liczby”, potrzebujemy
zmiennych, których wartościami są \emph{zbiory elementów}, a nie pojedyncze
elementy. Jest to główna motywacja logiki drugiego rzędu.

\begin{definicja}[Zmienne i formuły drugiego rzędu]\label{def:second-order-syntax}
Ustalmy język pierwszego rzędu $\mathcal L$ i jego zmienne
indywidualne $x,y,z,\ldots$. Rozszerzamy alfabet o nieskończenie
wiele \emph{zmiennych relacyjnych} $X^{(n)},Y^{(n)},\ldots$ dla
każdej dodatniej liczby naturalnej $n$ (jej \emph{arności}).
Jeśli $X^{(n)}$ ma arność $n$, a $t_1,\ldots,t_n$ są termami, to
$X^{(n)}(t_1,\ldots,t_n)$ jest formułą atomową. Oprócz spójników
i kwantyfikatorów pierwszego rzędu dopuszczamy
$\forall X^{(n)}\,\varphi$ i $\exists X^{(n)}\,\varphi$.
Formuły otrzymujemy przez skończone iterowanie tych konstrukcji;
nic innego nie uznajemy za formułę. Zmienną relacyjną nazywamy
\emph{wolną}, jeśli jej wystąpienie nie pozostaje w zasięgu
wiążącego ją kwantyfikatora.
Dla czytelności zmienną jednoargumentową zapisujemy często po
prostu jako $X$, a $X(x)$ odczytujemy „$x$ należy do zbioru $X$”.
\end{definicja}

\begin{przyklad}[Relacja a zbiór]
Napis $X(x)$ nie jest tym samym co $x=X$: pierwszy jest formułą
mówiącą, że element $x$ należy do zbioru reprezentowanego przez $X$;
drugi nie jest w opisanej składni poprawnym atomem, bo $x$ i $X$
są zmiennymi różnych rodzajów. Zapis $\exists X\,\forall x\,X(x)$
mówi, że istnieje zbiór zawierający każdy element uniwersum.
\end{przyklad}

\subsection{Semantyka pełna i semantyka Henkina}
\begin{definicja}[Spełnianie w pełnej semantyce drugiego rzędu]
\label{def:second-order-full-semantics}
Niech $\mathfrak A$ będzie strukturą pierwszego rzędu o
niepustym uniwersum $A$. Wartościowanie $s$ przypisuje zmiennym
indywidualnym elementy $A$, a zmiennej relacyjnej $X^{(n)}$
relację $s(X^{(n)})\subseteq A^n$. Dla atomu przyjmujemy
\[
 \mathfrak A,s\models X^{(n)}(t_1,\ldots,t_n)
 \quad\Longleftrightarrow\quad
 (t_1^{\mathfrak A}[s],\ldots,t_n^{\mathfrak A}[s])
 \in s(X^{(n)}).
\]
Pozostałe atomy, spójniki i kwantyfikatory indywidualne
interpretujemy jak w Definicji~\ref{def:sem-satisfaction}.
Dla kwantyfikatorów nowego rodzaju określamy:
\[
\begin{aligned}
\mathfrak A,s\models\forall X^{(n)}\,\varphi
 &\iff \text{dla każdego }D\subseteq A^n\text{ zachodzi }
 \mathfrak A,s[X^{(n)}\mapsto D]\models\varphi,\\
\mathfrak A,s\models\exists X^{(n)}\,\varphi
 &\iff \text{dla pewnego }D\subseteq A^n\text{ zachodzi }
 \mathfrak A,s[X^{(n)}\mapsto D]\models\varphi.
\end{aligned}
\]
Ta interpretacja jest \emph{pełną semantyką drugiego rzędu}:
zmienne relacyjne obejmują \emph{wszystkie} relacje odpowiedniej
arności na $A$.
\end{definicja}

\begin{definicja}[Semantyka Henkina]\label{def:second-order-henkin-semantics}
Niech $\mathfrak A$ będzie strukturą pierwszego rzędu
o niepustym uniwersum $A$. W \emph{semantyce Henkina} ustalamy
niepustą rodzinę $\mathcal S_n\subseteq\mathcal P(A^n)$ dla
każdej rozważanej arności $n$. Wartościowanie zmiennych
indywidualnych przyjmuje wartości w $A$, natomiast
wartościowanie każdej zmiennej relacyjnej $X^{(n)}$
musi spełniać $s(X^{(n)})\in\mathcal S_n$,
również gdy $X^{(n)}$ jest zmienną wolną.
Kwantyfikatory $\forall X^{(n)}$ i $\exists X^{(n)}$
przebiegają wówczas tylko po $D\in\mathcal S_n$.
Atom $X^{(n)}(\vec t)$ ma tę samą interpretację co powyżej. Nie zakładamy, że
$\mathcal S_n=\mathcal P(A^n)$; ewentualne dodatkowe warunki
na rodziny $\mathcal S_n$ (na przykład aksjomaty
\emph{rozumienia}) należy wysłowić oddzielnie.
\end{definicja}

\begin{przyklad}[Jedno zdanie, różne zakresy kwantyfikatora]
Niech $A=\N$ i $S(n)=n+1$. Formuła
\[
\forall X\,\Bigl(\bigl(X(0)\land
\forall x\,(X(x)\to X(S(x)))\bigr)
\to\forall x\,X(x)\Bigr)
\tag{$\mathrm{Ind}_2$}
\]
wyraża indukcję \emph{dla każdego zbioru} $X$.
Przy pełnej semantyce $X$ przebiega cały zbiór potęgowy
$\mathcal P(\N)$: założenie oznacza, że $X$ zawiera zero
i jest zamknięty na następnik, a wniosek, że $X=\N$.
W semantyce Henkina $X$ przebiega tylko rodzinę
$\mathcal S_1$. Wówczas formuła może być spełniona,
nawet jeżeli istnieją inne zewnętrzne podzbiory $A$
niedostępne dla kwantyfikatora.
\end{przyklad}

\begin{przyklad}[Model, w którym dwa odczytania indukcji się różnią]
Niech $A=\Z$, $0^{\mathfrak A}=0$ i
$S^{\mathfrak A}(z)=z+1$. Przy pełnej semantyce rozważmy
podzbiór $D=\{z\in\Z:z\geq0\}$. Zawiera on $0$
i jest zamknięty na funkcję $S^{\mathfrak A}$,
lecz nie zawiera $-1$. Podstawiając $D$ za $X$,
sprawdzamy, że poprzednik implikacji $\mathrm{Ind}_2$
jest prawdziwy, a jej następnik fałszywy.
Tak więc $\mathfrak A\not\models\mathrm{Ind}_2$
w pełnej semantyce.

Ustalmy teraz semantykę Henkina z
$\mathcal S_1=\{\Z\}$ (a dla pozostałych arności
niepustymi rodzinami, np. $\mathcal S_n=\{\Z^n\}$).
Jedyną wartością zmiennej jednoargumentowej $X$
jest wtedy całe $\Z$. Dla $X=\Z$ implikacja
$\mathrm{Ind}_2$ jest prawdziwa: spełniony jest zarówno
jej poprzednik, jak i następnik.
Zatem zdanie indukcji drugiego rzędu jest spełnione
w tej strukturze Henkina. Model nie spełnia jednak
odpowiedniego warunku pełnej semantyki.
Ten przykład nie twierdzi, że otrzymana struktura
spełnia wszystkie pozostałe aksjomaty arytmetyki Peana;
ilustruje wyłącznie rolę \emph{zakresu kwantyfikatora $X$}.
\end{przyklad}

\begin{uwaga}[Dlaczego to rozróżnienie jest ważne?]
Twierdzenia Gödla o pełności oraz Löwenheima--Skolema
w naszych dalszych rozdziałach dotyczą \emph{logiki pierwszego
rzędu}. Nie wolno przenosić ich bez dodatkowych założeń
na pełną semantykę drugiego rzędu. W
późniejszym twierdzeniu o kategoryczności arytmetyki Peana
zastosujemy właśnie formułę $\mathrm{Ind}_2$ i wykażemy, dlaczego
pełna semantyka zapewnia kategoryczność arytmetyki Peana
(nie ma niestandardowych modeli spełniających pełną indukcję
drugiego rzędu). Kategoryczność oznacza tu, że każde dwa modele
są izomorficzne, a nie że istnieje algorytm rozstrzygający
wszystkie zdania teorii.
\end{uwaga}

\section{Dlaczego pełność nie wynika z samej definicji?}
Zdefiniowaliśmy prawdziwość w strukturze. Nie rozstrzyga to,
czy każde zdanie prawdziwe we wszystkich modelach ma formalny dowód.
W kolejnym rozdziale określimy reguły dowodzenia, a następnie
wykażemy ich \emph{poprawność}: dowód formalny gwarantuje
wynikanie semantyczne. Odwrotność, czyli \emph{pełność Gödla},
wymaga odrębnej konstrukcji.